% Standalone appendix document - compiles on its own:
%   latexmk -pdf -outdir=build 06-meeting-minutes-2026-10-28.tex
\documentclass[a4paper,11pt]{article}
\usepackage[danish]{babel}
\usepackage{../../appendix-style}

\title{13.6 Mødereferat -- 28-10}
\author{SW4PRJ4 -- Group 3}
\date{28-10-2026}

\begin{document}
\maketitle

\section*{Mødeinformation}

\textbf{Dato:} 28-10-2026 (onsdag)

\textbf{Tid:} 10:00

\textbf{Sted:} Incuba

\textbf{Referent:} Frederik

\section*{Deltagere}

\begin{itemize}
    \item Alexander H
    \item Alexander D
    \item Frederik
    \item Ida D
    \item Ida F
    \item Malthe
    \item Marie
    \item Oskar
\end{itemize}

\textbf{Fraværende:}

\begin{itemize}
    \item Ingen
\end{itemize}

\section*{Dagsorden}

\begin{enumerate}
    \item Sprint review -- sprint 3
    \item Retrospective
    \item Sprint planning -- sprint 4
    \item Næste vejledermøde
    \item Eventuelt
\end{enumerate}

\section*{Noter}
\subsection*{1. Sprint review}

\begin{itemize}
    \item Velocity for sprint 3: committed 31,5 point, done 24,5 point, spillover 7 point, scope 91,5 point.
    \item PMU-70, PMU-74, PMU-75, PMU-76, PMU-81, PMU-82 og PMU-85 er merget og flyttes til \emph{Done}.
    \item PMU-77, PMU-78, PMU-79 og PMU-80 er ikke færdige og går videre til sprint 4.
    \item Kaneo er opdateret, så det, vi ikke nåede, står på boardet.
    \item Alle kort i \emph{In Review} er godkendt.
    \item Vi viste ikke en demo i dag. Den bliver i stedet vist til næste sprint review, når de resterende kort er færdige, og diagrammerne er lavet.
\end{itemize}

\subsection*{2. Retrospective}

\begin{itemize}
    \item \textbf{Opfølgning:} Vi har arbejdet test først, og det fungerede rigtig godt.
    \item \textbf{Mange PR'er samme dag:} Det skyldes, at vi kun arbejder sammen om onsdagen, og derfor bliver det meste merget på én dag.
    \item \textbf{Rapportkortene:} Det er hovedsageligt Frederik og Oskar, der har haft travlt i sprinten, og derfor er rapportkortene ikke nået. Der bliver fokuseret mere på rapporten i sprint 4.
    \item \textbf{Hvad gik godt?}
    \begin{itemize}
        \item Vi har fået en bedre forståelse af Keycloak.
        \item Gruppefordelingen fungerer godt. Undergrupperne har en god størrelse, og opgaverne er godt spredt ud.
    \end{itemize}
    \item \textbf{Hvad gik skidt?}
    \begin{itemize}
        \item Lokalet i Incuba ligger langt fra kaffemaskinen.
        \item Vi opdagede, at controllerne kaldte repository-klasserne direkte og dermed ikke gik gennem services, som arkitekturen kræver. Det er rettet i PMU-87.
        \item Integrationstestene er skrevet uden en egentlig plan, fordi vi endnu ikke har haft om integrationstest i undervisningen.
    \end{itemize}
    \item \textbf{Hvad ændrer vi?}
    \begin{itemize}
        \item Vi bruger dependency trees og det øvrige fra undervisningen i softwaretest, når vi skriver tests fremover.
        \item Hver person flytter selv sine kort på boardet, når de skifter status.
    \end{itemize}
\end{itemize}

\subsection*{3. Sprint planning}

\begin{itemize}
    \item Sprint 4 bliver en sprint på én uge (28/10 -- 3/11), så vi kan få de sidste ting fra sprint 3 gjort færdige.
    \item \textbf{Sprint goal:} De ufærdige kort fra sprint 3 gøres færdige, og rapporten bringes op på niveau med koden med diagrammer og gennemskrevne dokumenter.
    \item \textbf{Committed:}
    \item Indhold i sprinten:
    \begin{itemize}
        \item Diagrammer generelt, fx flowdiagrammer, sekvensdiagrammer og pakkediagram.
        \item Dokumenterne gennemgås, og der skrives mere i dem.
        \item De kort fra sprint 3, vi ikke nåede.
        \item Databasen skal på plads inden mødet med vejlederen tirsdag.
    \end{itemize}
\end{itemize}

\subsection*{4. Næste vejledermøde}

\begin{itemize}
    \item Næste vejledermøde er tirsdag 3/11.
    \item Spørgsmål til vejlederen:
    \begin{itemize}
        \item Hvilke diagrammer forventes der i rapporten ud over flow-, sekvens- og pakkediagrammer?
        \item Vil han anbefale mere kode i en sprint på én uge, eller skal den bruges på rapporten?
        \item Hvad anbefaler han, at vi går i gang med i næste sprint?
    \end{itemize}
\end{itemize}

\section*{Beslutninger}

\begin{itemize}
    \item Spillover-kortene PMU-77 -- PMU-80 beholder label \emph{Sprint 3} og får \emph{Sprint 4} oveni.
    \item Sprint 4 forkortes til én uge (28/10 -- 3/11), så de resterende kort fra sprint 3 og rapportarbejdet kan blive færdigt, inden vi starter på ny funktionalitet.
    \item Demoen udskydes til næste sprint review, så den kan vise det færdige flow sammen med diagrammerne.
\end{itemize}

\section*{Opgaver}

\begin{itemize}
    \item \textbf{Ansvarlig:} Alle \textbf{Opgave:} Flyt kort på boardet, når de skifter status \textbf{Frist:} løbende
    \item \textbf{Ansvarlig:} Alle \textbf{Opgave:} Database på plads til vejledermødet \textbf{Frist:} tirsdag 3/11
\end{itemize}

\section*{Næste møde}

\begin{itemize}
    \item Tirsdag 3/11: vejledermøde.
    \item Onsdag 4/11: sprint review med demo, retrospective og sprint planning for sprint 5.
\end{itemize}

\end{document}
